\documentstyle[% 


righttag% 


,11pt% 


,amssymb% 


,russian%               remove in Eng. 


,izvru%                 Set izven% in Eng. 


% ,izvru-ks             for brief communications 


]{amsart} 


 


% 


 


\newcommand{\const}{\operatorname{const}} 


\newcommand{\myH}{\mathrm H} %% seek for another symbol 


\newcommand{\eqdef}{\overset{\mathrm{def}}{=}{}} 


\newcommand{\tts}[1]{} 


 


 


%\hoffset=-15mm%                  For Eng. 


\setcounter{page}{3} 


\god{2004} 


\nomer{9~(508)} %                        Remove in Eng. 


%\nomer{Vol.\,48, No.\,9}%               Set in Eng. 


%\str{\pageref{begin}--\pageref{end}}%  Set in Eng. 


\udk{514.756} 


 


\author{\it Д.А.\,Абруков} 


% NB:   ^^ remove in Eng. 


\title{Полнота фундаментального объекта  


поверхности,\\ не принадлежащей абсолюту\\


проективно-метрического пространства} 


 


\date{} 


%\pagestyle{myheadings}%         Set in Eng. 


\pagestyle{plain} %              Remove in Eng. 


\begin{document} 


%\thispagestyle{plain}%          Set in Eng. 


\maketitle 


\label{begin} 


 


 


 


В работе исследуется распределение $m$-мерных линейных элементов 


$\Im$ ($m<n-1$), вложенное в проективно-метрическое пространство 


$K_n$ с абсолютом $Q_{n-1}$. Доказано, что в дифференциальной окрестности 


первого порядка распределение $\Im$ порождает инвариантно присоединенное 


к нему гиперполосное распределение $m$-мерных линейных 


элементов $\myH$ [1], для которого данное распределение и является базисным. 


Изучается $m$-мерная поверхность $V_m$ ($m<n-1$), текущая точка которой 


не принадлежит абсолюту $Q_{n-1}$ проективно-метрического пространства 


$K_n$. Показано, что во второй дифференциальной окрестности текущей точки 


поверхности $V_m$ индуцируется ассоциированная с ней гиперполоса $H(V_m)$. 


Доказано, что фундаментальный геометрический объект пятого порядка 


поверхности $V_m\subset K_n$ является полным. 


 


Индексы принимают следующие значения: 


\begin{gather*} 


\ov I,\ov K,\ov L=\ov{0,n};\quad I,K,L,P,Q=\ov{1,n};\\ 


i,j,k,p,q,s,l,t=\ov{1,m};\quad u,v,w,z,y=\ov{m+1,n-1};\quad 


\alpha,\beta,\gamma=\ov{m+1,m}. 


\end{gather*} 


 


Операция внешнего дифференцирования обозначена буквой $D$, 


внешнего умножения --- символом $\wedge$; $\omega^{\ov I}_{\ov K}$ и 


$\Omega^{\ov I}_{\ov K}$ --- дифференциальные 


формы Пфаффа, характеризующие инфинитезимальные перемещения подвижного 


репера. Оператор $\nabla$ действует по следующему закону: 


$$ 


\nabla K^\alpha_{in}=dK^\alpha_{in}-K^\alpha_{tn}\omega^t_i- 


K^\alpha_{in}\omega^n_n+K^\beta_{in}\omega^\alpha_\beta, 


$$ 


при фиксированных главных параметрах этот оператор обозначается через 


$\nabla_\delta$, формы $\omega^{\ov I}_{\ov K}$ --- через $\pi^{\ov I}_{\ov 


K}$. Оператор $\wt\nabla$ действует по закону 


$$ 


\wt\nabla T^\alpha_{in}=dT^\alpha_{in}-T^\alpha_{tn}\Omega^t_i- 


T^\alpha_{in}\Omega^n_n+T^\beta_{in}\Omega^\alpha_\beta. 


$$ 


По индексам, заключенным в круглые скобки, производится операция 


циклирования: 


$$ 


a_{(ij)}=a_{ij}+a_{ji}. 


$$ 


 


{\bf 1.} Рассмотрим $n$-мерное проективное пространство $P_n$, отнесенное к 


подвижному реперу $R'=\{B_{\ov K}\}$; деривационные формулы репера $R'$ и 


уравнение структуры проективного пространства имеют соответственно вид [2]: 


\begin{gather} 


dB_{\ov I}=\Omega^{\ov L}_{\ov I}B_{\ov L},\tag{1${}_{\mathrm a}$} \\ 


D\Omega^{\ov I}_{\ov K}=\Omega^{\ov L}_{\ov K}\wedge 


\Omega^{\ov I}_{\ov L},\quad 


\Omega^{\ov L}_{\ov L}=0. \tag{1${}_{\mathrm b}$} 


\end{gather} 


\addtocounter{equation}{1}





\notag{}Известно [3], [4], что проективно-метрическим пространством $K_n$ называется 


пространство $P_n$, в котором задана неподвижная гиперквадрика $Q_{n-1}$ 


(абсолют): 


\begin{equation} 


G_{\ov I\,\ov K}y^{\ov I}y^{\ov K}=0,\quad 


G_{\ov I\,\ov K}=G_{\ov K\,\ov I}, 


\end{equation} 


где $y^{\ov I}$ --- координаты точек $M'\in Q_{n-1}$. Согласно работе [5] 


условием неподвижности гиперквадрики (2) является выполнение 


дифференциальных уравнений 


\begin{equation} 


dG_{\ov I\,\ov K}-G_{\ov I\,\ov L}\Omega^{\ov L}_{\ov K}- 


G_{\ov L\,\ov K}\Omega^{\ov L}_{\ov I}= \Omega G_{\ov I\,\ov K}, 


\end{equation} 


где $\Omega$ --- некоторая форма Пфаффа. 


 


Известно [4], что при $B_0\notin Q_{n-1}$ за счет определенной нормировки 


вершин репера $R'$ становятся главными формы 


\begin{equation} 


\Omega^0_0=-\frac{G_{0L}}{G_{00}}\Omega^L_0. 


\end{equation} 


В силу соотношения (4) и за счет нормировки коэффициентов $G_{\ov I\,\ov K}$ 


гиперквадрики можно уравнение (2) абсолюта $Q_{n-1}$ и условие его 


неподвижности (3) записать соответственно в виде 


\begin{gather} 


c_{IK}y^I y^K+\frac{1}{G}(G_{I0}y^I+G y^0)^2=0,\notag\\ 


% 


dc_{IK}-c_{IL}\Omega^L_K-c_{LK}\Omega^L_I=-\frac{1}{G} 


(c_{IL}G_{K0}+c_{KL}G_{I0})\Omega^L_0, \tag{5${}_{\mathrm a}$} \\ 


% 


dG_{I0}-G_{L0}\Omega^L_I-G\Omega^0_I= 


c_{IL}\Omega^L_0,\tag{5${}_{\mathrm b}$} 


\end{gather} 


где 


$$ 


c_{IK}=G_{IK}-\frac{G_{I0}G_{K0}}{G},\quad c_{IK}=c_{KI},\quad 


G=G_{00}=\const\ne0. 


$$ 


В силу выражения (4) деривационные уравнения $(1_{\mathrm a})$ в 


проективно-метрическом пространстве $K_n$ запишутся в виде 


$$ 


dB_0=\bigg(-\frac{G_{0L}}{G}B_0+B_L \bigg)\Omega^L_0,\quad 


dB_I=\Omega^{\ov L}_I B_{\ov L}. 


$$ 


Уравнения структуры $(1_{\mathrm b})$ проективно-метрического пространства 


$K_n$ примут вид 


\begin{gather} 


D\Omega^I_0=-\frac{G_{0L}}{G}\Omega^L_0\wedge\Omega^I_0+ 


\Omega^L_0\wedge\Omega^I_L,\quad D\Omega^0_I=\Omega^L_I\wedge 


\Omega^0_L-\frac{G_{0L}}{G}\Omega^0_I 


\wedge\Omega^L_0, \tag{6${}_{\mathrm a}$} \\ 


% 


D\Omega^0_0=\Omega^L_0\wedge\Omega^0_L,\quad 


D\Omega^I_K=\Omega^0_K\wedge\Omega^I_0+\Omega^L_K\wedge 


\Omega^I_L,\quad \Omega^L_L=\frac{G_{0L}}{G}\Omega^L_0. 


\tag{6${}_{\mathrm b}$} 


\end{gather} 


\addtocounter{equation}{2}


 


{\bf 2.} Как и в пространстве проективной связности, уравнения распределения 


$m$-мерных линейных элементов $\Im$ в проективном-метрическом 


пространстве $K_n$ в репере нулевого порядка $R'=\{B_{\ov K}\}$ (точка 


$B_0\notin Q_{n-1}$ совпадает с центром 


распределения $\Im$, а точки $B_i$ принадлежат текущей плоскости 


$\Pi_m(B_0)$ распределения $\Im$) записываются в виде [6]: 


\begin{equation} 


\Omega^\alpha_i=M^\alpha_{iK}\Omega^K_0. 


\end{equation} 


 


Продолжая уравнения (7), с учетом соотношения (4) имеем 


\begin{gather} 


\wt\nabla M^\alpha_{ij}=M^\alpha_{ijK}\Omega^K_0,\\ 


% 


\wt\nabla M^\alpha_{i\beta}-M^\alpha_{ik}\Omega^k_\beta- 


\delta^\alpha_\beta\Omega^0_i=M^\alpha_{i\beta K} 


\Omega^K_0,\notag 


\end{gather} 


где 


$$ 


M^\alpha_{i[jk]}=M^\alpha_{i\beta}M^\beta_{[jk]}+\frac{1}{G} 


M{}^\alpha_{i[j}G{}_{k]0},\quad M^\alpha_{ik\beta}- 


M^\alpha_{i\beta k}=M^\alpha_{i\gamma}M^\gamma_{k\beta}+ 


\frac{2}{G}M{}^\alpha_{i[k}G{}_{\beta]0}. 


$$ 


Совокупность функций $\{M^\alpha_{ij}\}$ образует тензор первого порядка 


(вообще говоря, несимметричный). 


 


С учетом уравнений (7) условия (5) неподвижности абсолюта $Q_{n-1}$ 


запишутся в виде 


\begin{gather} 


dG_{i0}-G_{j0}\Omega^j_i-G\Omega^0_i=(c_{iK}+G_{\alpha0} 


M^\alpha_{iK})\Omega^K_0, \tag{9${}_{\mathrm a}$} \\ 


% 


dG_{\alpha0}-G_{\beta0}\Omega^\beta_\alpha-G_{j0} 


\Omega^j_\alpha-G\Omega^0_\alpha= 


c_{\alpha K}\Omega^K_0, \tag{9${}_{\mathrm b}$} \\ 


% 


\wt\nabla c_{ij}=\bigg(c{}_{\alpha(i}M{}^\alpha_{j)K} 


-\frac{1}{G}c{}_{K(i}G{}_{j)0} \bigg)\Omega^K_0,\tag{9${}_{\mathrm c}$}\\ 


% 


\wt\nabla c_{i\alpha}-c_{is}\Omega^s_\alpha=\bigg(c_{\alpha\beta} 


M^\beta_{iK}-\frac{1}{G}c{}_{K(i}G{}_{\alpha)0} \bigg) 


\Omega^K_0, \tag{9${}_{\mathrm d}$} \\ 


\wt\nabla c_{\alpha\beta}-c{}_{s(\alpha}\Omega{}^s_{\beta)}= 


-\frac{1}{G}c{}_{K(\alpha}G{}_{\beta)0}\Omega^K_0. 


\tag{9${}_{\mathrm e}$} 


\end{gather} 


\addtocounter{equation}{1}





\noindent{}Совокупность функций $\{c_{ij}\}$ образует симметричный тензор; предполагая 


его невырожденным, т.\,е. $c=|c_{ij}|\ne0$, имеем поле взаимного тензора 


$c^{ij}$: 


\begin{equation} 


c^{ik}c_{kj}=\delta^i_j,\quad \wt\nabla c^{ij}= 


c^{is}c^{tj}\bigg[\frac{1}{G}G{}_{0(t}c{}_{s)K}- 


c{}_{\alpha(s}M{}^\alpha_{t)K}\bigg]\Omega^K_0. 


\end{equation} 


 


Рассмотрим охваты 


\begin{gather} 


A_{\alpha\beta}\eqdef c_{\alpha\beta}-c^{ij}c_{i\alpha} 


c_{j\beta},\quad A_{[\alpha\beta]}=0;\\ 


b^\alpha\eqdef c^{ij}M^\alpha_{ij},\quad 


b_\alpha\eqdef A_{\alpha\beta}b^\beta; 


\end{gather} 


их компоненты в силу (8), $(9_{\mathrm{d,e}})$, (10) удовлетворяют 


дифференциальным уравнениям 


\begin{equation} 


\begin{array}{c} 


\wt\nabla A_{\alpha\beta}=A_{\alpha\beta K}\Omega^K_0;\quad 


\wt\nabla b^\alpha=b^\alpha_K\Omega^K_0;\quad 


\wt \nabla b_\alpha=b_{\alpha K}\Omega^K_0,\\[2mm] 


%


\wt\nabla b_{\alpha K}-\delta^\beta_K b_\beta\Omega^0_\alpha- 


b_\beta M^\beta_{iK}\Omega^i_\alpha=b_{\alpha KL}\Omega^L_0. 


\end{array} 


\end{equation} 


Следовательно, все три охвата (11), (12) образуют тензор; в частности 


$b_\alpha$ и $b^\alpha$ суть тензоры первого порядка. 


 


Предположим, что $m<n-1$, т.\,е. распределение $m$-мерных линейных 


элементов $\Im$ не является распределением гиперплоскостных элементов. 


Инвариантная гиперплоскость $\xi_0$, уравнение которой в репере нулевого 


порядка имеет вид $b_\alpha y^\alpha=0$, содержит текущий $m$-мерный линейный 


элемент $\Pi_m$ распределения $\Im$. Следовательно, распределение 


$m$-мерных линейных элементов $\Im$ в первой дифференциальной окрестности 


внутренним образом порождает инвариантно присоединенное к нему 


гиперполосное распределение $m$-мерных линейных элементов $\myH$ [7], для 


которого данное распределение является базисным. 


 


Предположим, $B_n\notin\xi_0$, что равносильно условию $b_n\ne0$. 


Очевидно, точки $A_u=B_u-\frac{b_u}{b_n}B_n$ находятся в общем положении и 


лежат в гиперплоскости $\xi_0$, т.\,е. $\xi_0=[B_0,B_i,A_u]$. 


 


Инфинитезимальные перемещения репера $R=\{A_{\ov K}\}$ первого порядка 


($A_0=B_0$, $A_i=B_i$, $A_u=B_u-\frac{b_u}{b_n}B_n$, $A_n=B_n$) определяются 


уравнениями 


\begin{equation} 


dA_{\ov I}=\omega^{\ov L}_{\ov I}A_{\ov L}. 


\end{equation} 


 


Сравнивая $(1_{\mathrm a})$ и (14), с использованием уравнений (7), (13) 


находим 


\begin{gather} 


\displaybreak[3]


\omega^0_0=\Omega^0_0,\quad \omega^i_0=\Omega^i_0,\quad 


\omega^u_0=\Omega^u_0,\quad \omega^n_0=\Omega^n_0- 


\frac{b_u}{b_n}\Omega^u_0,\notag \\ 


% 


\omega^0_j=\Omega^0_j,\quad \omega^i_j=\Omega^i_j,\quad 


\omega^u_j=\Omega^u_j,\quad \omega^n_j=\frac{1}{b_n} 


b_\alpha\Omega^\alpha_j,\notag \\ 


% 


\omega^0_u=\Omega^0_u-\frac{b_u}{b_n}\Omega^0_n,\quad 


\omega^i_u=\Omega^i_u-\frac{b_u}{b_n}\Omega^i_n,\quad 


\omega^v_u=\Omega^v_u-\frac{b_u}{b_n}\Omega^v_n,\\ 


% 


\omega^n_u=\frac{1}{(b_n)^2}(b_u b_{nK}-b_n b_{uK}) 


\Omega^K_0,\notag \\ 


% 


\omega^0_n=\Omega^0_n,\quad \omega^i_n=\Omega^i_n,\quad 


\omega^u_n=\Omega^u_n,\quad \omega^n_n=\Omega^n_n 


+\frac{b_u}{b_n}\Omega^u_n.\notag 


\end{gather} 


 


Уравнения структуры проективно-метрического пространства $K_n$ в 


репере $R=\{A_{\ov K}\}$ первого порядка имеют вид 


$$ 


D\omega^{\ov I}_{\ov K}=\omega^{\ov L}_{\ov K}\wedge 


\omega^{\ov I}_{\ov L},\quad \omega^{\ov L}_{\ov L}=0. 


$$ 


 


Известно [7], что относительно репера $R=\{A_{\ov K}\}$ уравнения 


гиперполосного распределения $m$-мерных линейных элементов $\myH$, 


ассоциированного 


с распределением $\Im$, запишутся в виде 


\begin{equation} 


\begin{array}{c} 


\omega^n_i=\Lambda^n_{iK}\omega^K_0,\quad 


\omega^v_i=\Lambda^v_{iK}\omega^K_0,\\[2mm] 


\omega^n_u=A^n_{uK}\omega^K_0. 


\end{array} 


\end{equation} 


Из соотношений (7), (15), (16) находим 


\begin{equation} 


\Lambda^n_{iK}=\frac{1}{b_n}b_\alpha M^\alpha_{iK},\quad 


\Lambda^u_{iK}=M^u_{iK},\quad A^n_{uK}=\frac{1}{(b_n)^2} 


(b_u b_{nK}-b_n b_{uK}). 


\end{equation} 


Продолжая уравнения (16), в частности, имеем 


\begin{gather} 


\nabla A^n_{uj}+A^n_{uj}\omega^0_0-\Lambda^n_{sj} 


\omega^s_u=A^n_{ujK}\omega^K_0, \tag{18${}_a$} \\ 


\nabla\Lambda^n_{ij}+\Lambda^n_{ij}\omega^0_0= 


\Lambda^n_{ijK}\omega^K_0. \tag{18${}_b$} 


\end{gather} 


\addtocounter{equation}{1}


 


Отметим, что в силу соотношений (17) функции $\Lambda^n_{iK}$, 


$\Lambda^v_{iK}$ относятся к 


первой дифференциальной окрестности, а функции $A^n_{uK}$ --- ко второй 


дифференциальной окрестности текущего элемента распределения $\Im$. 


 


При условии невырожденности тензора $\Lambda^n_{ij}$ (см. $(18_b)$)


$\Lambda=|\Lambda^n_{ij}|\ne0$, в силу 


уравнений $(18_a)$ согласно лемме Н.М.\,Остиану [8] возможна частичная 


канонизация репера $R=\{A_{\ov K}\}$, при которой функции $A^n_{uj}$ можно 


привести к нулю, при этом становятся главными формы 


$\omega^i_u=N^i_{uK}\omega^K_0$. В полученном 


репере второго порядка уравнения гиперполосного распределения 


$m$-мерных линейных элементов $\myH$, ассоциированного с распределением 


$\Im$, запишутся в виде 


\begin{equation} 


\omega^n_i=\Lambda^n_{iK}\omega^K_0,\quad \omega^v_i=\Lambda^v_{iK} 


\omega^K_0,\quad \omega^n_u=A^n_{u\alpha}\omega^\alpha_0,\quad 


\omega^i_u=N^i_{uK}\omega^K_0. 


\end{equation} 


Здесь функции $\Lambda^n_{iK}$, $\Lambda^v_{iK}$ относятся к 


дифференциальной окрестности первого порядка, функции $A^n_{u\alpha}$ --- к 


окрестности второго порядка, а функции 


$N^i_{vj}$ --- к дифференциальной окрестности третьего порядка текущего 


элемента распределения $m$-мерных линейных элементов $\Im$. 


 


Проведенная специализация репера $R=\{A_{\ov K}\}$ имеет следующую 


геометрическую характеристику [7]: 


$$ 


\Lambda=|\Lambda^n_{ij}|\ne0\ \Leftrightarrow \ \Pi_m 


\cap \Pi_\rho\equiv A_0, 


$$ 


где $\Pi_\rho$ --- характеристика гиперплоскости $\xi_0\equiv\Pi_{n-1}$ при 


смещениях центра $A_0$ вдоль кривых $l$, принадлежащих [6] базисному 


распределению, причем $A_u\in\Pi_\rho$, т.\,е. $\rho=n-m-1$. 


 


Итак, справедлива 


 


\begin{thm} 


Распределение $m$-мерных линейных элементов $\Im$ 


$(m<n-1)$ в дифференциальной окрестности первого порядка порождает 


инвариантно присоединенное к нему гиперполосное распределение 


$m$-мерных линейных элементов $\myH$, для которого данное распределение 


$\Im$ является базисным. Распределение $\myH$ регулярно тогда и только 


тогда, когда невырожден тензор $b_\alpha M^\alpha_{ij}$ первого порядка. 


Уравнения регулярного гиперполосного распределения $\myH$ в репере второго 


порядка $R=\{A_{\ov K}\}$ имеют вид $(19)$. 


\end{thm} 


 


 


{\bf 3.} Вернемся к подвижному реперу $R'=\{B_{\ov K}\}$. Деривационные 


формулы репера $R'$ и уравнения структуры проективно-метрического 


пространства $K_n$ в репере $R'$ имеют соответственно вид (6). 


 


Рассмотрим $m$-мерную поверхность $V_m$ проективно-метрического 


пространства $K_n$, текущая точка $B_0$ которой не принадлежит абсолюту 


$Q_{n-1}$ (см. (5)). Как и в проективном пространстве $P_n$, в репере 


$R'=\{B_{\ov K}\}$ первого порядка ($B_0\in V_m$, $B_i\in T_m(B_0)$) 


уравнения поверхности $V_m$ в $K_n$ записываются в виде [9] 


\begin{equation} 


\Omega^\alpha_0=0. 


\end{equation} 


 


Трехкратное продолжение уравнений (20) с учетом соотношения (4) 


приводит к следующим дифференциальным уравнениям компонент полей 


фундаментальных объектов второго $\{M^\alpha_{ij}\}$ и третьего 


$\{M^\alpha_{ij},M^\alpha_{ijk}\}$ порядков $V_m\subset K_n$: 


\begin{gather} 


\Omega^\alpha_i=M^\alpha_{ij}\Omega^j_0,\quad 


M^\alpha_{[ij]}=0,\notag \\ 


\wt\nabla M^\alpha_{ij}=M^\alpha_{ijk}\Omega^k_0,\quad 


M^\alpha_{i[jk]}=\frac{1}{G}M{}^\alpha_{i[j}G{}_{k]0},\quad 


M^\alpha_{[ij]k}=0,\notag \\ 


\wt\nabla M^\alpha_{ijk}+M{}^\alpha_{k(i}\Omega{}^0_{j)}- 


M{}^\beta_{(ij}M{}^\alpha_{k)s}\Omega^s_\beta= 


M^\alpha_{ijkl}\Omega^l_0,\quad M^\alpha_{ij[kl]}= 


\frac{1}{G}M{}^\alpha_{ij[k}G{}_{l]0}. 


\end{gather} 


Совокупность функций $\{M^\alpha_{ij}\}$ образует симметричный тензор второго 


порядка, $\{M^\alpha_{ij},M^\alpha_{ijk}\}$ --- фундаментальный геометрический 


объект третьего порядка поверхности $V_m$. 


 


Из уравнений (21) следует, что в дифференциальной окрестности точки 


$B_0\in V_m$ ниже третьего порядка невозможно построить внутреннюю 


инвариантную нормализацию поверхности $V_m\subset K_n$. Ниже будет 


показано, что такую нормализацию возможно построить в третьей 


дифференциальной окрестности. 


 


Отметим, что согласно работе Остиану [10] имеет место аналогия 


между геометрией $m$-мер\-ной поверхности проективно-метрического 


пространства $K_n$ и геометрией распределения $m$-мер\-ных линейных элементов 


(см. п.\,2) пространства $K_n$ в той части, которая определяется 


последовательностью полей фундаментальных подобъектов 


$\{M^\alpha_{ij}\},\{M^\alpha_{ij},M^\alpha_{ijk_1}\},\dotsc, 


\{M^\alpha_{ij},M^\alpha_{ijk_1},\dotsc$, $M^\alpha_{ijk_1\dotsc k_s}\}$ 


%%%%%%%%%%%%% join                     ^^^^


распределения $\Im$. 


 


В силу последнего утверждения уравнения неподвижности абсолюта 


$Q_{n-1}$ (см. (5)) и уравнения тензоров $c^{ij}$, $A_{\alpha\beta}$, 


$b^\alpha$, $b_\alpha$ в репере $R'=\{B_{\ov K}\}$ 


первого порядка будут иметь вид (9), (10), (13) соответственно, с той лишь 


разницей, что для поверхности $V_m\subset K_n$ дополнительно выполняются еще 


и уравнения (20). Отметим, что на поверхности $V_m\subset K_n$ ($m<n-1$) 


$b_\alpha$ и $b^\alpha$ являются тензорами второго порядка. 


 


Пусть $m<n-1$, т.\,е. поверхность $V_m$ в $K_n$ отлична от гиперповерхности. 


Инвариантная гиперплоскость $\xi_0$, уравнение которой в репере 


$R'=\{B_{\ov K}\}$ первого порядка имеет вид $b_\alpha y^\alpha=0$, в каждой 


точке $B_0\in V_m$ содержит касательную плоскость $T_m(B_0)$ поверхности 


$V_m$. Следовательно, $m$-мерная поверхность $V_m\subset K_n$ во второй 


дифференциальной окрестности внутренним образом порождает инвариантно 


присоединенную к ней гиперполосу $H(V_m)$, для которой данная поверхность 


является базисной [11]. 


 


Аналогично п.\,2 предположим, что $B_n\notin\xi_0$; последнее равносильно 


тому, что $b_n\ne0$. Очевидно, точки $A_u=B_u-\frac{b_u}{b_n}B_n$ находятся в 


общем положении и лежат в гиперплоскости $\xi_0$, где $\xi_0=[B_0,B_i,A_u]$. 


 


Инфинитезимальные перемещения репера $R=\{A_{\ov K}\}$ (см. п.\,2) второго 


порядка ($A_0=B_0$, $A_i=B_i$, $A_u=B_u-\frac{b_u}{b_n}B_n$, $A_n=B_n$) 


определяются уравнениями (14). 


 


Связь между формами $\Omega^{\ov K}_{\ov I}$ и $\omega^{\ov K}_{\ov I}$ 


имеет вид (15). Здесь следует учесть, что для поверхности $V_m\subset K_n$ 


справедливы уравнения (20). 


 


Известно [1], [12], что относительно репера $R=\{A_{\ov K}\}$ второго порядка 


уравнения гиперполосы $H(V_m)$, ассоциированной с поверхностью 


$V_m$, запишутся в виде 


$$ 


\omega^\alpha_0=0,\quad \omega^n_i=\Lambda^n_{ij}\omega^j_0,\quad 


\omega^v_i=\Lambda^v_{ij}\omega^j_0,\quad \omega^n_u= 


\Lambda^n_{uj}\omega^j_0. 


$$ 


 


Аналогично п.\,2 в предположении $\Lambda=|\Lambda^n_{ij}|\ne0$ возможна 


частичная канонизация репера $R=\{A_{\ov K}\}$, имеющая следующую 


геометрическую характеристику [1]: вершины $A_\alpha$ выбраны таким образом, 


чтобы $(n-m)$-мерная плоскость $[A_0,A_\alpha]$ пересекала гиперплоскость 


$\xi_0$ по ее характеристике $\Pi_{n-m-1}(A_0)$, т.\,е. 


$A_v\in\Pi_{n-m-1}(A_0)$. В полученном репере третьего порядка $R=\{A_{\ov 


K}\}$ уравнения гиперполосы $H(V_M)$, ассоциированной 


с поверхностью $V_m$, примут вид [1], [13] 


\begin{equation} 


\begin{array}{c} 


\omega^\alpha_0=\omega^n_v=0,\quad \omega^n_i=\Lambda^n_{ij} 


\omega^j_0,\quad \omega^v_i=\Lambda^v_{ij}\omega^j_0,\\[2mm] 


\omega^i_u=N^i_{uj}\omega^j_0,\quad \Lambda^n_{[ij]}= 


\Lambda^u_{[ij]}=0,\quad \Lambda{}^n_{s[i} 


N{}^s_{|v|j]}=0, 


\end{array} 


\end{equation} 


где 


$$ 


\Lambda^n_{ij}=\frac{1}{b_n}b_\alpha M^\alpha_{ij},\quad 


\Lambda^u_{ij}=M^u_{ij},\quad \Lambda^n_{us}=\frac{1}{(b_n)^2} 


(b_u b_{ns}-b_n b_{us}). 


$$ 


Здесь функции $\Lambda^n_{ij}$, $\Lambda^v_{ij}$ относятся к дифференциальной 


окрестности второго порядка, а функции $N^i_{vj}$ --- к дифференциальной 


окрестности четвертого порядка текущей точки поверхности $V_m$. 


 


\begin{thm} 


Поверхность $V_m\subset K_n$ $(m<n-1)$ в дифференциальной окрестности 


второго порядка порождает инвариантно присоединенную к 


ней гиперполосу $H(V_m)$, для которой данная поверхность является базисной. 


Гиперполоса $H(V_m)$ регулярна тогда и только тогда, когда невырожден 


тензор $b_\alpha M^\alpha_{ij}$ второго порядка. Уравнения регулярной 


гиперполосы $H(V_m)$ в репере третьего порядка $R=\{A_{\ov K}\}$ имеют вид 


$(22)$. 


\end{thm} 


 


 


Уравнения (9) неподвижности абсолюта $Q_{n-1}$ проективно-метрического 


пространства $K_n$ с учетом соотношений (15) и уравнений 


(20), (22) перепишутся, в частности, в виде 


\begin{align} 


dG_{i0}-G_{j0}\omega^j_i-G\omega^0_i&=(\dotsc)_{is} 


\omega^s_0,\notag \\ 


\nabla c_{iu}-c_{iw}\bigg(\frac{b_u}{b_n} \bigg)\omega^w_n-c_{is} 


\bigg(\frac{b_u}{b_n} \bigg)\omega^s_n-c_{in} 


\Omega^n_u&=(\dotsc)_{ius}\omega^s_0,\\ 


\nabla c_{in}-\bigg\{c_{iw}-c_{in}\bigg(\frac{b_w}{b_n} \bigg) 


\bigg\}\omega^w_n-c_{is}\omega^s_n&=(\dotsc)_{ins}\omega^s_0.\notag 


\end{align} 


 


Отметим, что выражение (4) с использованием (15), (20) запишется в 


виде 


\begin{equation} 


\omega^0_0=-\frac{G_{0s}}{G}\omega^s_0. 


\end{equation} 


 


Продолжая уравнения (22), с учетом выражений (24) имеем 


\begin{gather} 


\nabla\Lambda^n_{ij}=\Lambda^n_{ijk}\omega^k_0,\quad 


\Lambda^n_{i[jk]}=\frac{1}{G}\Lambda{}^n_{i[j}G{}_{k]0},\\ 


% 


\nabla\Lambda^u_{ij}+\Lambda^n_{ij}\omega_n^u=\Lambda^u_{ijk} 


\omega^k_0,\quad \Lambda^u_{i[jk]}=\frac{1}{G} 


\Lambda{}^u_{i[j}G{}_{k]0},\\ 


% 


\nabla N^i_{uj}-\delta^i_j\omega^0_u=N^i_{ujk}\omega^k_0,\quad 


N^i_{u[jk]}=\frac{1}{G}N{}^i_{u[j}G{}_{k]0},\\ 


% 


\nabla\Lambda^n_{ijk}+\Lambda{}^n_{k(i}\omega{}^0_{j)}- 


\Lambda{}^n_{(ij}\Lambda{}^n_{k)s}\omega^s_n= 


\Lambda^n_{ijks}\omega^s_0,\\ 


% 


\Lambda^n_{ij[ks]}=\Lambda^n_{il}\Lambda{}^w_{j[k} 


N{}^l_{|w|s]}+\Lambda^n_{jl}\Lambda{}^w_{i[k} 


N{^l_{|w|s]}}-\frac{2}{G}\Lambda{}^n_{ij[k}G{}_{s]0}.\notag 


\end{gather} 


В силу невырожденности тензора $\Lambda^n_{ij}:\Lambda=|\Lambda^n_{ij}|\ne0$ 


можно ввести обращенный тензор $\Lambda^{ij}_n$ второго порядка: 


\begin{equation} 


\Lambda^{ik}_n\Lambda^n_{kj}=\delta^i_j,\quad 


\nabla\Lambda^{ij}_n=-\Lambda^{is}_n\Lambda^{tj}_n 


\Lambda^n_{stl}\omega^l_0. 


\end{equation} 


Функция $\Lambda$ есть относительный инвариант второго порядка гиперполосы 


$H(V_m)$: 


$$ 


d\ln\Lambda-2\omega^k_k+m\omega^n_n=\Lambda_l\omega^l_0,\quad 


\Lambda_l=\Lambda^{ji}_n\Lambda^n_{ijl}. 


$$ 


Продолжая последнее уравнение, имеем 


\begin{equation} 


\nabla\Lambda_i+2\omega^0_i-(m+2)\Lambda^n_{ij}\omega^j_n= 


\Lambda_{ij}\omega^j_0,\quad \Lambda_{[ij]}= 


2\Lambda{}^u_{s[i}N{}^s_{|u|j]}+\frac{1}{G} 


\Lambda{}_{[i}G{}_{j]0}. 


\end{equation} 


 


В дифференциальных окрестностях второго и четвертого порядков 


строим поля квазитензоров $a^u_n$ и $a_u$: 


$$ 


a^u_n=\frac{1}{m}\Lambda^u_{ij}\Lambda^{ij}_n,\quad 


a_u=\frac{1}{m}N^i_{ui}, 


$$ 


компоненты которых в силу (25)--(27), (29) удовлетворяют 


дифференциальным уравнениям 


\begin{equation} 


\nabla a^u_n+\omega^u_n=a^u_{ns}\omega^s_0,\quad 


\nabla a_u-\omega^0_u=\wt a_{us}\omega^s_0. 


\end{equation} 


 


Уравнения (10) в силу (15), (20), (22) перепишутся в виде 


\begin{equation} 


\nabla c^{ij}=(\dotsc)^{ij}_s\omega^s_0. 


\end{equation} 


Используя уравнения (13), с учетом (15), (20), (22) имеем 


\begin{equation} 


d\bigg(\frac{b_u}{b_n} \bigg)-\bigg(\frac{b_w}{b_n} \bigg) 


\omega^w_u+\bigg(\frac{b_u}{b_n} \bigg)\omega^n_n- 


\bigg(\frac{b_u b_w}{(b_n)^2} \bigg)\omega^w_n-\Omega^n_u= 


(\dotsc)^n_{us}\omega^s_0. 


\end{equation} 


 


 


{\bf 4.} Регулярная гиперполоса $H(V_m)$ называется оснащенной в смысле 


Нордена--Чакмазяна [3], [14], если заданы два поля квазитензоров 


[7]: 


\begin{equation} 


\nabla\nu^i_n+\omega^i_n=\nu^i_{nj}\omega^j_0,\quad 


\nabla\nu_i+\omega^0_i=\nu_{ij}\omega^j_0, 


\end{equation} 


определяющих соответственно поля нормалей первого и второго родов на 


данном подмногообразии. 


 


Согласно уравнениям (23), (25), (29)--(33) в третьей дифференциальной 


окрестности внутренним образом определяются поля нормалей первого 


и второго родов гиперполосы $H(V_m)$: 


\begin{align} 


P^i_n&\eqdef-\frac{1}{m+2}\Lambda^{ik}_n 


\bigg(\Lambda_k+\frac{2G_{k0}}{G} \bigg),\tag{35${}_a$} \\ 


% 


P_i&\eqdef\frac{1}{2}\bigg(\Lambda_i-(m+2)\Lambda^n_{ij}c^{\,jk} 


\bigg[c_{kn}+\bigg\{c_{kw}-c_{kn}\bigg(\frac{b_w}{b_n} \bigg) 


\bigg\}a^w_n\bigg]\bigg).\tag{35${}_b$} 


\end{align} 


\addtocounter{equation}{1}


 


Итак, справедлива 


 


\begin{thm} 


Внутренним образом определенная инвариантная нормализация 


гиперполосы $H(V_m)$, а значит, и внутренняя инвариантная нормализация 


ее базисной поверхности $V_m\subset K_n$ $(m<n-1)$ определяются в 


третьей дифференциальной окрестности текущей точки $A_0\in V_m$ полями 


квазитензоров $(35)$. 


\end{thm} 


 


 


{\bf 5.} Говорят, что гиперполоса $H(V_m)$ оснащена в смысле Э.\,Картана 


[15], если в каждой точке $A_0$, принадлежащей ее базисной поверхности 


$V_m$, поставлена в соответствие плоскость $N_{n-m-1}(A_0)$ размерности 


$(n-m-1)$, не имеющая общих точек с касательной плоскостью $T_m(A_0)$ к 


поверхности $V_m$. Плоскость $N_{n-m-1}(A_0)$ в каждой точке $A_0\in V_m$ 


можно задать точками 


$$ 


M_u=\nu_u A_0+A_u;\quad M_n=\nu_n A_0+\nu^i_n A_i+\nu^u_n A_u. 


$$ 


Функции, входящие в последние соотношения, удовлетворяют уравнениям 


[1], [7]


\begin{gather*} 


\nabla\nu_u+\omega^0_u=\nu_{uj}\omega^j_0,\quad 


\nabla\nu_n+\nu^i_n\omega^0_i+\nu^u_n\omega^0_u+ 


\omega^0_n=\nu_{nj}\omega^j_0,\\ 


\nabla\nu^u_n+\omega^u_n=\nu^u_{nj}\omega^j_0,\quad 


\nabla\nu^i_n+\omega^i_n=\nu^i_{nj}\omega^j_0. 


\end{gather*} 


 


Итак, оснащение в смысле Э.\,Картана гиперполосы $H(V_m)$ равносильно 


заданию полей геометрических объектов $\{\nu^i_n\}$, 


$\{\nu^i_n,\nu^u_n,\nu_n\}$, $\{\nu_u\}$, $\{\nu^u_n\}$. 


 


Из уравнений (25)--(28) следует, что внутреннее инвариантное оснащение 


в смысле Э.\,Картана гиперполосы $H(V_M)$ в $K_n$ ($m<n-1$), а 


следовательно, и ее базисной поверхности $V_m$ невозможно построить в 


дифференциальной окрестности порядка ниже четвертого. Покажем, что такое 


оснащение внутренним образом возможно построить в четвертой 


дифференциальной окрестности точки $A_0\in V_m$. 


 


Продолжая уравнения (34), с учетом (24) имеем 


$$ 


\nabla_\delta\nu^i_{nj}+\Lambda^n_{kj}\nu{}^{(k}_n\pi{}^{i)}_n 


-\delta^i_j(\nu^k_n\pi^0_k+\pi^0_n)-N^i_{uj}\pi^u_n=0. 


$$ 


 


В качестве оснащающих полей геометрических объектов $\{\nu^i_n\}$, 


$\{\nu^u_n\}$, $\{\nu^i_n,\nu^u_n,\nu_n\}$, $\{\nu_u\}$ 


по аналогии с [1], [7] берем 


\begin{equation} 


\nu^u_n\eqdef a^u_n,\quad \nu_u\eqdef -a_u,\quad 


\nu_n\eqdef\frac{1}{m}(\nu^j_{nj}-\Lambda^n_{kj} 


\nu^k_n\nu^j_n)-a^u_n a_u, 


\end{equation} 


а в качестве функций $\{\nu^i_n\}$ --- квазитензор $P^i_n$ третьего порядка 


$(35_a)$. 


 


Таким образом, имеет место 


 


\begin{thm} 


Внутреннее инвариантное оснащение в смысле Э.\,Картана 


гиперполосы $H(V_m)$ в $K_n$ $($а значит, и ее базисной поверхности 


$V_m\subset K_n$ $(m>n-1))$ строится в четвертой дифференциальной окрестности 


текущей точки $A_0$ поверхности $V_m$ полями объектов $(36)$, где 


$\nu^i_n=P^i_n$. 


\end{thm} 


 


Из теорем 3 и 4 непосредственно следует 


 


\begin{thm} 


Фундаментальный геометрический объект не ниже пятого 


порядка поверхности $V_m$, не принадлежащей абсолюту 


проективно-метрического пространства $K_n$ $(n<m-1)$, является полным. 


\end{thm} 


 


 


{\bf 6.} Известно [1], [12], что порядок полного внутреннего 


фундаментального объекта поверхности $V_m$ проективного пространства $P_n$ 


($2<m<n-1$) не превосходит шести. В [1], [12] при доказательстве 


данного утверждения имеется в виду, что главный фундаментальный 


тензор $\Lambda^n_{ij}$ гиперполосы $H(V_m)$, ассоциированной с 


поверхностью $V_m\subset P_n$ ($2<m<n-1$), определяется в третьей 


дифференциальной окрестности текущей точки $A_0\in V_m$; из [16] следует, 


что порядок этого геометрического объекта не ниже пяти. 


 


Отметим, что в случае поверхности частного класса, а именно, поверхности 


Картана $V_m\subset P_{2m}$, $m>2$, доказано [17], что ее фундаментальный 


геометрический объект пятого порядка является полным. 


 


Так как для гиперполосы $H(V_m)$, ассоциированной с поверхностью 


$V_m\subset K_n$ ($m<n-1$), ее главный фундаментальный тензор 


$\Lambda^n_{ij}$ относится ко второй дифференциальной окрестности точки 


$A_0\in V_m$ (см. (28)), то, следуя схеме доказательства полноты 


фундаментального геометрического объекта регулярной гиперполосы 


(доказательство приведено, напр., в [1]), заключаем, что 


справедлива 


 


\begin{thm} 


Порядок полного внутреннего фундаментального объекта 


поверхности $V_m$, не принадлежащей абсолюту проективно-метрического 


пространства $K_n$ $(n<m-1)$, не превосходит пяти. 


\end{thm} 


 


В силу теорем 5 и 6 справедлива 


 


\begin{thm} 


Порядок полного внутреннего фундаментального объекта поверхности $V_m$, не 


принадлежащей абсолюту проективно-метрического пространства $K_n$ 


$(n<m-1)$, равен пяти. 


\end{thm} 


 


В случае регулярности гиперполосы $H(V_m)$ в $K_n$, ассоциированной с 


поверхностью $V_m\subset K_n$ ($m<n-1$), справедливы результаты, полученные 


в [1] для гиперполосы, погруженной в проективное пространство 


$P_n$. Перечислим основные из них (применительно к поверхности $V_m\subset 


K_n$). 


 


1. Инвариантное присоединение регулярной гиперполосы $H(V_m)$ к 


поверхности $V_m\subset K_n$ позволяет построить двойственную теорию [1] 


многообразия $V_m$. 


 


2. В дифференциальной окрестности четвертого порядка в каждой 


точке $A_0$ поверхности $V_m\subset K_n$ имеем канонический пучок 


инвариантных нормалей первого рода. Из этого пучка выделяются, в частности, 


нормаль Фубини, директриса Вильчинского и т.\,д. 


 


Используя двойственную теорию гиперполосы, нетрудно получить 


поле пучка нормалей второго рода поверхности $V_m\subset K_n$. 


 


3. В дифференциальной окрестности четвертого порядка точки $A_0\in V_m$ 


имеем поле инвариантных соприкасающихся гиперквадрик $Q^2_{n-1}$ 


гиперполосы $H(V_m)$ (а следовательно, и ее базисной поверхности $V_m$), 


относительно которых касательная плоскость $T_m(A_0)$ и характеристика 


$\Pi_{n-m-1}(A_0)$ полярно сопряжены. 


 


Используя инвариантное условие квадратичности гиперполосы 


$H(V_m)$ в $P_n$ ($m<n-1$) (см. [1], [18]), нетрудно найти критерий 


принадлежности поверхности $V_m$ неподвижной гиперквадрике $Q^2_{n-1}$. 


 


4. Инвариантное присоединение регулярной гиперполосы $H(V_m)$ к 


$V_m\subset K_n$ позволяет изучать на поверхности $V_m$: 


\begin{itemize} 


\item[а)] две двойственные аффинные связности без кручения, 


индуцируемые нормализацией $H_m$ в смысле Нордена--Чакмазяна, а также 


индуцируемые при этом двойственные нормальные связности $\overset{1}{D}$ и 


$\overset{2}{D}$ и их подсвязности, также двойственные между собой (см. 


[1]); 


\item[б)] две двойственные проективные связности без кручения, 


индуцируемые [1] оснащением поверхности в смысле Картана. 


\end{itemize} 


 


 


\begin{thebibliography}{99} 


 


\bibitem{s6} %1


Столяров А.В. {\em Двойственная теория оснащенных многообразий}. 


-- Чебоксары, 1994. -- 290~с. 


\bibitem{fin} %2


Фиников С.П. {\em Метод внешних форм Картана в дифференциальной 


геометрии}. -- М.-Л.: ГИТТЛ, 1948. -- 432 с. 


\bibitem{nap} %3


Норден А.П. {\em Пространства аффинной связности}. -- М.: Наука, 


1976. -- 432 с. 


\bibitem{s7} %4


Столяров А.В. {\em Внутренняя геометрия проективно-метрического 


пространства} // Дифференц. геометрия многообразий фигур. -- Изд-во 


Калинингр. ун-та. -- 2001. -- Вып.\,32. -- С.\,94--101. 


\bibitem{lap1} %5


Лаптев Г.Ф. {\em Дифференциальная геометрия погруженных многообразий} // 


Тр. Матем. о-ва. -- М., 1953. -- Т.\,2. -- С.\,275--382. 


\bibitem{lapos} %6


Лаптев Г.Ф., Остиану Н.М. {\em Распределения $m$-мерных линейных 


элементов в пространстве проективной связности}. I // Тр. Геометрич. 


семин. ВИНИТИ, -- 1971. -- Т.\,3. -- С.\,49--94. 


\bibitem{s3} %7


Столяров А.В. {\em Проективно-дифференциальная геометрия регулярного 


гиперполосного распределения $m$-мерных линейных элементов} //


Итоги науки техники. Пробл. геометрии. -- М.: ВИНИТИ, 1975. --


T.\,7. -- C.\,117--151.


\bibitem{ost1} %8


Остиану Н.М. {\em О канонизации подвижного репера погруженного 


многообразия} // Rev. math. pures et appl. -- 1962. -- V.\,7. -- \No\,2. 


-- P.\,231--240. 


\bibitem{lap2} %9


Лаптев Г.Ф. {\em Дифференциальная геометрия многомерных поверхностей} 


// Итоги науки и техн. Геометрия. -- M.: ВИНИТИ,


1965. -- С.\,5--64. 


\bibitem{ost3} %10


Остиану Н.М. {\em Распределения $m$-мерных линейных элементов в 


пространстве проективной связности}. II // Тр. Геометрич. семин. 


ВИНИТИ. -- 1971. -- Т.\,3. -- С.\,95--114. 


\bibitem{vag} %11


Вагнер В.В. {\em Теория поля локальных гиперполос} // Тр. семин. по 


векторн. и тензорн. анализу. -- М: Изд-во МГУ. -- 1950. 


-- Вып.\,8. -- С.\,197--272. 


\bibitem{s5} %12


Столяров А.В. {\em Приложение теории регулярных гиперполос к изучению 


геометрии многомерных поверхностей проективного пространства} 


// Изв. вузов. Математика. -- 


1976. -- \No\,2. -- С.\,111--113. 


\bibitem{s1} %13


Столяров А.В. {\em О фундаментальных объектах регулярной гиперполосы} // 


Изв. вузов. Математика. -- 1975. -- \No\,10. -- С.\,97-99. 


\bibitem{cha} %14


Чакмазян А.В. {\em Двойственная нормализация} // Докл. АН 


АрмССР. -- 1959. -- Т.\,28. -- \No\,4. -- С.\,151--157. 


\bibitem{cart} %15


Cartan E. {\em Les espaces \'a connexion projective} // Тр. семин. по 


векторн. и тензорн. анализу. -- М.: Изд-во МГУ. -- 1937. -- 


Вып.\,4. -- С.\,147--159. 


\bibitem{ost2} %16


Остиану Н.М. {\em О геометрии многомерной поверхности проективного 


пространства} // Тр. Геометрич. семин. ВИНИТИ. 


-- 1966. -- Т.\,1. -- С.\,239--263. 


\bibitem{s4} %17


Столяров А.В. {\em О внутренней геометрии поверхности Картaна} // 


Дифференц. геометрия многообразий фигур. -- Изд-во Калинингр. ун-та. 


-- 1976. -- Вып.\,7. -- С.\,111--118. 


\bibitem{s2} %18


Столяров А.В. {\em Условие квадратичности регулярной гиперполосы} 


// Изв. вузов. Математика. -- 1975. - \No\,11. -- С.\,106--108. 


 


\end{thebibliography} 


 


\vskip 2cm 


 


\post{\it Чувашский государственный}{\it Поступила} 


\post{\it педагогический университет}{22.04.2002} 


 


\label{end} 


\end{document} 


 


